% Abhakseite „Das kann ich“ – Zone nur im Gesamt (\mitzone)
%% TODO Ich-kann-Sätze: Platzhalter wie in den Aufgabendateien
\begin{abhakseite}
\ifdefined\mitzone
\abhakgruppe{Kennst du schon}
\abhak{1}{Rechtwinkliges Dreieck erkennen, Rechtwinkelmarke (Bogen mit Punkt), Hypotenuse gegenüber dem rechten Winkel, Katheten am rechten Winkel}
\abhak{2}{Winkel benennen (α, β, γ, auch β1, β2, ε), Winkel messen und schätzen, spitz oder stumpf}
\abhak{3}{Taschenrechner}
\abhak{4}{Gleichung mit einem Bruch nach einer Größe umstellen, Größe im Zähler (mal) oder im Nenner (geteilt), mit Strich}
\abhak{5}{Bruch als Verhältnis lesen und in eine Dezimalzahl umwandeln}
\abhak{6}{Satz des Pythagoras als zweiter Weg zur dritten Seite, Vergleich der Ergebnisse}
\abhak{7}{Winkelsumme im Dreieck, die beiden spitzen Winkel des rechtwinkligen Dreiecks als Ergänzung, Winkel aus Teilwinkeln (Differenz), Strecke aus Teilstrecken}
\abhak{8}{Längeneinheiten vor dem Rechnen angleichen (m und km, cm und m), Ergebnis mit Einheit, Runden auf eine Dezimale}
\abhak{9}{Höhe, Diagonale, Überstand am Trapez, Fußpunkt auf der Verlängerung beim Parallelogramm, halbe Diagonale im Drachen}
\abhak{10}{Gleichung mit einem Bruch nach einer Größe umstellen, Größe im Zähler (mal) oder im Nenner (geteilt), mit Strich – Fehler finden}
\abhak{11}{Gleichung mit einem Bruch nach einer Größe umstellen, Größe im Zähler (mal) oder im Nenner (geteilt), mit Strich – selbst rechnen}
\fi
\abhakgruppe{Einheit 1 · Seite berechnen mit Sinus, Kosinus und Tangens}
\abhak{12}{sin, cos oder tan?}
\abhak{13}{Mal oder geteilt?}
\abhak{14}{Seite berechnen}
\abhak{15}{Seiten vom zweiten spitzen Winkel aus neu benennen (Gegen- und Ankathete tauschen)}
\abhak{16}{Seitenverhältnis aus gegebenen Längen als Dezimalzahl und Vergleich mit dem Taschenrechnerwert des Winkels}
\abhak{17}{Ergebnis prüfen: Kathete kürzer als Hypotenuse}
\abhak{18}{Seite berechnen – Fehler finden, Begründen, Anwendung}
\abhakgruppe{Einheit 2 · Winkel berechnen}
\abhak{19}{Winkel berechnen}
\abhak{20}{Winkel prüfen: die längere Kathete liegt dem größeren Winkel gegenüber}
\abhak{21}{Umkehrtaste am Taschenrechner (SHIFT sin, cos, tan; Anzeige sin⁻¹)}
\abhak{22}{Winkel berechnen – Fehler finden, Begründen, Anwendung}
\abhakgruppe{Einheit 3 · Rechtwinklige Teildreiecke in Figuren und Vermessung}
\abhak{23}{Teildreiecke und Vermessung}
\abhak{24}{Strecke aus Teilstrecken nach der Trigonometrie}
\abhak{25}{Teildreiecke und Vermessung – Fehler finden, Begründen, Anwendung}
\abhakgruppe{Einheit 4 · Sinussatz}
\abhak{26}{Sinussatz}
\abhak{27}{Sinussatz – Fehler finden, Begründen, Anwendung}
\end{abhakseite}
